% ============================================================
%  ACADEMIA ZENIT — GUÍA CLASE 1: ENSAYO PSU 2005
%  Matemática · 3° Medio
% ============================================================

\documentclass[11pt, letterpaper]{article}

% -- Paquetes de Compatibilidad y Corrección --
\usepackage[T1]{fontenc}
\usepackage[utf8]{inputenc}
\usepackage[spanish, es-tabla]{babel}

\usepackage[margin=2cm, top=2.2cm, bottom=2.2cm]{geometry}
\usepackage{amsmath, amssymb, amsthm}
\usepackage{mathtools}
\usepackage{array, booktabs, tabularx}
\usepackage{xcolor}
\usepackage{graphicx}
\usepackage{enumitem}
\usepackage{fancyhdr}
\usepackage{tcolorbox}
\usepackage{tikz}
\usepackage{pgfplots}
\pgfplotsset{compat=1.18}
\usepackage{lastpage}

% -- Colores Institucionales Academia Zenit --
\definecolor{zenitnavy}{HTML}{1A3A5C}
\definecolor{zenitblue}{HTML}{2E6DA4}
\definecolor{zenitgold}{HTML}{C9820A}
\definecolor{zenitlightgray}{HTML}{F5F5F5}
\definecolor{blocka}{HTML}{1A3A5C}
\definecolor{blockb}{HTML}{1A6B3C}

% -- Estilos de tcolorbox --
\tcbuselibrary{skins, breakable}
\tcbset{
  instrbox/.style={
    enhanced,
    colback=white,
    colframe=zenitgold,
    left=8pt, right=8pt, top=6pt, bottom=6pt,
    arc=2pt,
    leftrule=4pt,
  }
}

% -- Encabezado de bloque --
\newcommand{\bloqueheader}[3]{%
  \vspace{10pt}
  \noindent
  \begin{tikzpicture}[baseline]
    \fill[#2] (0,0) rectangle (2.5cm, 0.7cm);
    \node[white, font=\bfseries] at (1.25cm, 0.35cm) {Bloque #1};
  \end{tikzpicture}%
  \begin{tikzpicture}[baseline]
    \fill[#2!10!white] (0,0) rectangle (\linewidth - 2.6cm, 0.7cm);
    \node[color=#2, font=\bfseries\large, anchor=west] at (0.2cm, 0.35cm) {#3};
  \end{tikzpicture}%
  \vspace{4pt}
}

% -- Configuración de Encabezado --
\pagestyle{fancy}
\fancyhf{}
\fancyhead[L]{\small \textbf{Academia Zenit} $\cdot$ Ensayo PSU 2005}
\fancyhead[R]{\small Página \thepage\ de \pageref{LastPage}}
\fancyfoot[L]{\small Prohibida su reproducción total o parcial.}

% ============================================================
\begin{document}
% ============================================================

\begin{center}
  {\Huge\bfseries\color{zenitnavy} ACADEMIA ZENIT}\\[4pt]
  {\LARGE\bfseries\color{zenitblue} ENSAYO MATEMÁTICA — FACSÍMIL 2005}\\[8pt]
  \rule{\linewidth}{1pt}
\end{center}

\begin{tcolorbox}[instrbox]
  \textbf{Instrucciones:} Esta prueba consta de 70 preguntas. Las figuras no están necesariamente dibujadas a escala. Dispones de 2 horas y 15 minutos para responder.
\end{tcolorbox}

\vspace{6pt}

% ============================================================
% --- BLOQUE A: NÚMEROS Y PROPORCIONALIDAD ---
% ============================================================
\bloqueheader{A}{blocka}{Números y Proporcionalidad}

\begin{enumerate}[label=\textbf{\arabic*.}, leftmargin=*, itemsep=10pt]

% ----------------------------------------------------------------
\item
\begin{minipage}[t]{\linewidth}
Determine el resultado de la siguiente operación:
\[
  \dfrac{9}{8} - \dfrac{3}{5} =
\]
\\
\begin{minipage}{\linewidth}
\vspace{7mm}
\begin{tikzpicture}[every node/.style={anchor=north west}]

% --- Primer nodo: alternativas ---
\node (texto) at (0,0) {
\begin{minipage}{0.95\linewidth}
\begin{enumerate}[label=\Alph*), leftmargin=14pt]
  \item $0{,}15$
  \item $0{,}5$
  \item $0{,}52$
  \item $0{,}525$
  \item $2$
\end{enumerate}
\end{minipage}
};

% --- Segundo nodo: figura ---
\node (figura) at (7,0) {
};

\end{tikzpicture}
\end{minipage}
\vspace{7mm}
\end{minipage}

% ----------------------------------------------------------------
\setcounter{enumi}{3}
\item
\begin{minipage}[t]{\linewidth}
Calcule el valor de la expresión:
\[
  \dfrac{3^{-1} + 4^{-1}}{5^{-1}}
\]
\\
\begin{minipage}{\linewidth}
\vspace{7mm}
\begin{tikzpicture}[every node/.style={anchor=north west}]

% --- Primer nodo: alternativas ---
\node (texto) at (0,0) {
\begin{minipage}{0.95\linewidth}
\begin{enumerate}[label=\Alph*), leftmargin=14pt]
  \item $\dfrac{12}{35}$
  \item $\dfrac{35}{12}$
  \item $7$
  \item $5$
  \item $\dfrac{5}{7}$
\end{enumerate}
\end{minipage}
};

% --- Segundo nodo: figura ---
\node (figura) at (7,0) {
};

\end{tikzpicture}
\end{minipage}
\vspace{7mm}
\end{minipage}

\end{enumerate}

% ============================================================
% --- BLOQUE B: ÁLGEBRA Y FUNCIONES ---
% ============================================================
\bloqueheader{B}{blockb}{Álgebra y Funciones}

\begin{enumerate}[label=\textbf{\arabic*.}, leftmargin=*, itemsep=10pt, start=26]

% ----------------------------------------------------------------
\item
\begin{minipage}[t]{\linewidth}
En la figura, las rectas $L_1$ y $L_2$ son perpendiculares. ¿Cuál de las siguientes opciones representa la ecuación de la recta $L_1$?
\\
\begin{minipage}{\linewidth}
\vspace{7mm}
\begin{tikzpicture}[every node/.style={anchor=north west}]

% --- Primer nodo: alternativas ---
\node (texto) at (0,0) {
\begin{minipage}{0.95\linewidth}
\begin{enumerate}[label=\Alph*), leftmargin=14pt]
  \item $y = \dfrac{5}{4}x - 2$
  \item $y = \dfrac{5}{4}(x - 2)$
  \item $y = \dfrac{4}{5}(x - 2)$
  \item $y = \dfrac{4}{5}x - 2$
  \item $y = -\dfrac{5}{4}(x - 2)$
\end{enumerate}
\end{minipage}
};

% --- Segundo nodo: figura ---
\node (figura) at (7,0) {
\begin{tikzpicture}[scale=0.7]
  \draw[-latex, thick] (-1,0) -- (6,0) node[right] {$x$};
  \draw[-latex, thick] (0,-1) -- (0,5) node[above] {$y$};
  \draw[thick, zenitnavy] (-0.5,4) -- (5.5,-0.8) node[above right] {$L_2$};
  \draw[thick, zenitblue] (1,-1) -- (5,4) node[right] {$L_1$};
  \node[left] at (0,4) {$4$};
  \node[below] at (5,0) {$5$};
  \node[below left] at (2,0) {$2$};
  \draw (2.4,0.5) -- (2.7,0.5) -- (2.7,0.15);
\end{tikzpicture}
};

\end{tikzpicture}
\end{minipage}
\vspace{7mm}
\end{minipage}

\end{enumerate}

\end{document}